% ============================================================================
%  Subdocumento 2. Comprensión del problema y de la necesidad
%  audIT Soluciones Tecnológicas SpA - Licitación TFEP-01/2026
% ============================================================================

\begin{resumenApertura}
Este subdocumento analiza los antecedentes del Caso 10 sobre la operación de Transportes Curimón S.A. y sus necesidades de información y coordinación. Describe la composición de la flota y de los conductores, las condiciones de conectividad y estacionalidad informadas, los actores del caso y los supuestos que podrían afectar el proyecto. Las cifras que requieren mayor detalle se identifican para cotejo con las bases y sus fuentes.
\begin{recibe}
  \item Análisis de la cadena de valor en seis fases operacionales y síntesis de brechas descritas en el Caso 10.
  \item Análisis de la estacionalidad frutícola informada (65\% de la demanda anual de frío en cinco meses) y de los cierres cordilleranos reportados para Los Libertadores; la autonomía local del sistema debe cubrir al menos las 72 horas exigidas por RT-03.10.
  \item Mapeo estratégico de poder e influencia sobre los 13 actores del ecosistema y criterios de arbitraje para las seis tensiones operacionales.
  \item Matriz de supuestos auténticos de ingeniería de proyectos con evaluación de impacto y planes de contingencia técnica.
  \item Cuatro anexos técnicos complementarios: catálogo de 25 requerimientos preliminares, inventario de 374 tractocamiones y 454 choferes, matriz de restricciones legales y 13 fichas pormenorizadas de stakeholders.
\end{recibe}
\end{resumenApertura}

Este subdocumento sintetiza los antecedentes del Caso 10 sobre la operación de Transportes Curimón S.A. en el marco de la Licitación Pública N.° TFEP-01/2026. Organiza las brechas de información, coordinación y control descritas en el caso, identifica las restricciones operacionales pertinentes y explicita los supuestos que requieren confirmación. No presenta estas cifras como resultado de un levantamiento independiente.

Este capítulo se articula orgánicamente con la totalidad de los subdocumentos y anexos de la propuesta técnica:
\begin{enumerate}
  \item Provee el fundamento fáctico e ingenieril para el diseño del alcance y la descomposición modular expuestos en el Subdocumento 3 (\textit{Esquema de Solución y Alcance}).
  \item Determina los requerimientos no funcionales de resiliencia, volumetría y desempeño que gobiernan la Arquitectura Lógica y Física desarrollada en el Subdocumento 4.
  \item Delimita los dominios de información, volumetría histórica y requerimientos de gobernanza formalizados en el Subdocumento 5 (\textit{Modelo y Gestión de Datos}).
  \item Define el contexto operacional bajo el cual se estructuran las metodologías de trabajo, los paquetes de la Estructura de Descomposición del Trabajo (EDT) y la matriz de riesgos analizados en los Subdocumentos 6, 7 y 8.
  \item Establece los umbrales basales de calidad, niveles de servicio y soporte continuado a 36 meses detallados en los Subdocumentos 9, 10 y 11.
  \item Se vincula de forma directa con la asignación de roles y perfiles del Subdocumento 12, la justificación del catálogo de innovaciones del Subdocumento 13 y la demostración de beneficios cuantificados en el Subdocumento 14.
\end{enumerate}

Asimismo, los inventarios detallados de requerimientos preliminares, el desglose pormenorizado del parque vehicular y conductores, la matriz exhaustiva de restricciones legales y las fichas completas de caracterización de actores se formalizan en el anexo técnico independiente \textit{AUDIT-Subdocumento2-Anexos.pdf}.

\section{Resumen Ejecutivo del problema}\label{sec:2-resumen-ejecutivo}

El diagnóstico de Transportes Curimón S.A. muestra una brecha entre las responsabilidades que la compañía asume ante sus mandantes y el control operacional que ejerce sobre los recursos que prestan el servicio de transporte interurbano. Según los antecedentes del caso, Curimón responde por la carga, el cumplimiento de los servicios y la observancia de las obligaciones aplicables. A la vez, el 60,4\% de la flota (226 de 374 tractocamiones) y el 56,8\% de los conductores asignables (258 externos y 196 propios) pertenecen a recursos subcontratados de 148 transportistas independientes, sobre los que Curimón no ejerce subordinación laboral directa.

Esta distribución de control plantea necesidades de coordinación y verificación entre Curimón y sus transportistas. La operación comprende 96.000 viajes anuales, 2,4 millones de toneladas y 41 millones de kilómetros recorridos. Se extiende por más de 3.000 kilómetros entre Antofagasta y Puerto Montt e incluye alrededor de 1.900 cruces por el Paso Los Libertadores. La asignación se apoya principalmente en telefonía, planillas y la Torre de Programación de San Bernardo. En ese contexto, el 26\% del kilometraje anual corresponde a recorridos de retorno sin carga (10,66 millones de kilómetros).

El caso también identifica diferencias de rentabilidad entre contratos. Tres de los ocho clientes principales operan bajo su costo técnico y representan en conjunto el 31\% de los despachos e ingresos. En uno de ellos se informa un margen negativo de 14\% durante cuatro años consecutivos. El prorrateo de costos según ingresos dificulta distinguir el resultado de cada ruta y contrato.

El caso señala además que el 71\% de los cobros por sobreestadía en instalaciones de clientes es objetado. La falta de registros objetivos de llegada, espera y salida limita la posibilidad de respaldar esos cobros y analizar sus causas.

En gestión documental y telemetría se observan brechas de registro y consolidación. Cerca de 6.000 vencimientos de licencias, permisos, revisiones técnicas, certificados y seguros se administran en cuatro planillas. El caso informa que no se han descargado registros de tacógrafos; 61 tractocamiones propios disponen de telemetría CAN J1939 aún no integrada; 34 camiones de terceros no tienen GPS; y las otras 340 unidades se distribuyen entre tres plataformas que no ofrecen una vista operacional unificada.

El ciclo de renovación previsto para 2029 plantea un condicionante comercial: el cliente exportador principal, que representa el 19\% de los ingresos según el caso, ha establecido cuatro requisitos para la continuidad del contrato:
\begin{enumerate}
  \item Trazabilidad integral y posicionamiento de carga en tiempo real para el 100\% de los despachos asignados.
  \item Emisión y tramitación de documentación electrónica de transporte sin redigitación manual ni soporte físico en papel (e-Docs para un volumen proyectado de 128.000 documentos anuales).
  \item Acreditación técnica fehaciente del cumplimiento de los límites de jornada laboral y descansos obligatorios del Artículo 25 bis del Código del Trabajo para la totalidad de los conductores (tanto propios como externos de terceros) en cada viaje asignado.
  \item Reportabilidad periódica y auditada de emisiones de gases de efecto invernadero (GEI / CO2e) por tonelada-kilómetro transportada, calculada bajo el estándar internacional del marco GLEC (\textit{Global Logistics Emissions Council}) / ISO 14083:2023.
\end{enumerate}

La falta de cumplimiento de esos requisitos pondría en riesgo la renovación del contrato asociado al 19\% de los ingresos indicado en el caso. Por ello, este subdocumento delimita las brechas operacionales, documentales y de información que deben considerarse en la propuesta, sin anticipar el diseño de la solución desarrollado en los capítulos posteriores.

\section{Comprensión del problema y de la necesidad}\label{sec:2-comprension}

Para comprender la raíz del problema operacional de Transportes Curimón S.A., es indispensable contextualizar la actividad en la cadena de valor del transporte terrestre de carga interurbana en Chile, examinando las restricciones geográficas, comerciales y los marcos legales que gobiernan la circulación por carretera.

En la \figref{2-flujo-operacional} se describe el flujo operacional y la cadena de valor característica del transporte de carga en Curimón, ilustrando la secuencia de procesos desde la recepción de la orden de transporte hasta la liquidación final y cierre de costos.

\figura{figuras/02-problema/Flujo_Operacional.png}{Flujo operacional y cadena de valor del transporte de carga en Curimón}{2-flujo-operacional}

El análisis de la \figref{2-flujo-operacional} muestra los principales quiebres de información y coordinación en las fases de la cadena de valor de Curimón:

\begin{enumerate}
  \item \textbf{Fase 1: Comercial y Recepción de Demanda:} Las solicitudes de 84 clientes se registran en el TMS 2013, por correo electrónico y mediante llamadas telefónicas. Al ingresar la orden no se valida automáticamente la compatibilidad de carga, como los requisitos de refrigeración, sustancias peligrosas o peso por eje. El caso también informa contratos con margen operacional negativo de hasta 14\%.
  \item \textbf{Fase 2: Programación y Asignación de Recursos:} Casamiento de órdenes de transporte con unidades de tracción (374 camiones) y tripulaciones (454 conductores). Los 22 despachadores de la Torre de Control de San Bernardo asignan viajes mediante llamadas telefónicas y consulta visual de cuatro planillas Excel desvinculadas entre sí. No existe validación algorítmica previa de jornada laboral (Artículo 25 bis del Código del Trabajo), lo que deriva en el despacho involuntario de choferes sin descanso suficiente, multiplicando el riesgo de siniestros viales por somnolencia.
  \item \textbf{Fase 3: Carga y Despacho en Instalaciones de Origen:} Posicionamiento del equipo en los recintos de clientes ($\approx 1.400$ puntos), faena de estiba y emisión de guías de despacho. Las demoras en andenes de carga promedian 3 horas 10 minutos y superan las 8 horas en la temporada agrícola de exportación. Al no existir integración telemática ni comprobación de peso por eje a la salida, los camiones quedan expuestos a infracciones por sobrepeso en plazas de pesaje del MOP (\caso{RT-04}{31}).
  \item \textbf{Fase 4: Tránsito en Ruta y Logística Interurbana:} La operación cubre la Ruta 5 entre Antofagasta y Puerto Montt e incluye cerca de 1.900 cruces anuales por el Paso Los Libertadores. Hay tramos con más de 80 kilómetros sin cobertura celular y cierres cordilleranos prolongados. En esas condiciones, la falta de transmisión y de registros de tacógrafo reduce la visibilidad sobre los vehículos y la jornada.
  \item \textbf{Fase 5: Arribo, Descarga y Cierre de Entrega:} Recepción física en destino, descarga de mercancías y firma del comprobante de entrega en papel (\textit{Proof of Delivery} [POD]). La falta de registro electrónico objetivo de entrada y salida genera la objeción del 71\% de los cobros de sobreestadías. Asimismo, el 4,2\% de las guías en papel se extravían o deterioran, y el 26\% de los kilómetros anuales (10,66 millones de km) se ejecutan retornando en vacío por falta de algoritmos de triangulación en tiempo real.
  \item \textbf{Fase 6: Liquidación, Costeo y Facturación:} Consolidación administrativa de fletes, pago a 148 transportistas terceros e imputación contable. Demanda ocho (8) profesionales durante nueve (9) días hábiles al cierre de cada mes, exhibiendo una tasa de error del 11\% en liquidaciones. El diésel se costea con un desfase de hasta 40 días, impidiendo detectar a tiempo desvíos de rendimiento de hasta un 19\% entre vehículos gemelos.
\end{enumerate}

\subsection{Diagnóstico holístico e integración causal de la cadena operacional}\label{sec:2-diagnostico}

Los antecedentes de Transportes Curimón S.A. muestran factores organizacionales y técnicos relacionados entre sí. La siguiente síntesis los agrupa por sus efectos sobre la operación:

\begin{enumerate}
  \item \textbf{Rentabilidad por contrato y asignación de costos:} El caso señala que tres contratos operan bajo costo y representan el 31\% de la actividad e ingresos corporativos; uno registra un margen negativo de 14\% durante cuatro años. La distribución de costos según ingresos dificulta atribuir combustible, peajes y tarifas de terceros a cada viaje y contrato, lo que limita la comparación de rentabilidad entre rutas.
  \item \textbf{Registros de tacógrafo y jornada de conductores externos:} El caso informa que no existen descargas históricas de tacógrafos y que Curimón no dispone de registros de las actividades previas de los 258 conductores subcontratados. Tras un volcamiento ocurrido en el kilómetro 312 de la Ruta 5 Sur a las 04:40 h, la investigación identificó que el conductor había prestado servicios para otra empresa antes del accidente y no habría contado con el descanso mínimo. El caso señala que el incidente derivó en la suspensión de contratos durante seis semanas y en una investigación de la autoridad laboral.
  \item \textbf{Consumo de combustible y rezago de información:} El combustible representa el 14\% de la estructura de costos. El caso informa una diferencia de rendimiento de hasta 19\% entre tractocamiones comparables y un retraso de hasta 40 días en la recepción de facturas. Este desfase demora el análisis de consumo y la identificación de sus causas.
  \item \textbf{Conectividad en alta montaña y zonas desérticas:} Los cierres del Paso Los Libertadores pueden prolongar la desconexión y existen tramos de más de 80 kilómetros sin cobertura celular en corredores del norte. Estas condiciones limitan la actualización de posición e itinerarios y la transmisión de datos durante el trayecto; el equipo embarcado debe conservar localmente la información por al menos 72 horas.
\end{enumerate}

\subsection{Marco normativo y técnico a considerar}\label{sec:2-marco-normativo}

Las siguientes normas y referencias técnicas enmarcan el diagnóstico. La aplicabilidad concreta de las obligaciones laborales y de seguridad debe confirmarse para cada situación operativa:

\begin{itemize}
  \item \textbf{Ley N.° 20.123 sobre trabajo en régimen de subcontratación:} La aplicación de las responsabilidades de la empresa principal depende de los supuestos y condiciones previstos en la normativa laboral. En el caso, los conductores externos dependen de 148 transportistas independientes. La necesidad operacional consiste en coordinar la verificación de antecedentes y obligaciones aplicables antes de asignar viajes; el alcance jurídico de los deberes de Curimón debe ser revisado por un especialista (\parencite{ley20123}).
  \item \textbf{Artículo 25 bis del Código del Trabajo:} A la fecha de la oferta (5 de octubre de 2026), el régimen de conductores de carga terrestre interurbana considera una jornada ordinaria de hasta 180 horas mensuales distribuida en no menos de 21 días, un descanso mínimo ininterrumpido de ocho horas dentro de cada período de 24 horas, y un máximo de cinco horas continuas de conducción seguido por un descanso mínimo de dos horas. La reforma de la Ley N.° 21.561 que sustituye el límite mensual por una de las alternativas de jornada previstas para el artículo 25 bis tiene vigencia diferida al 26 de abril de 2028. Los controles de jornada deberán considerar la normativa aplicable al momento de operar y validarse con asesoría laboral (\parencite{codigo_trabajo_25bis,res_ex_1213_2009}).
  \item \textbf{D.S. N.° 298/1994 (MTT) frente a D.S. N.° 43/2015 (MINSAL) para Sustancias Peligrosas (SUSPEL):} Curimón dispone de 18 unidades especializadas en transporte de sustancias químicas. El marco regulatorio impone dos normativas complementarias pero con ámbitos físicos diferenciados:
  \begin{itemize}
    \item El \textbf{D.S. N.° 298/1994} del MTT reglamenta el transporte de cargas peligrosas por calles y caminos; los requisitos aplicables al vehículo, la carga y el conductor deben cotejarse con el texto vigente (\parencite{ds298}).
    \item El \textbf{D.S. N.° 43/2015} del MINSAL regula el almacenamiento de sustancias peligrosas en instalaciones fijas. La aplicabilidad de sus exigencias a cada terminal, patio y actividad descritos en el caso debe determinarse según las condiciones del recinto y revisarse con el responsable de prevención de riesgos (\parencite{ds43minsal}).
  \end{itemize}
  \item \textbf{Ley de Tránsito N.° 18.290 y Ley N.° 21.377:} La normativa restringe la manipulación de teléfonos y otros dispositivos electrónicos durante la conducción, con condiciones y excepciones que deben comprobarse en el texto vigente y su reglamentación. El proyecto puede proponer controles de interfaz para reducir distracciones, pero el umbral de bloqueo, el uso de alertas de voz y las excepciones deben definirse mediante revisión normativa y pruebas de seguridad (\parencite{ley21377}).
  \item \textbf{Protección de datos personales:} A la fecha de esta oferta se debe considerar la Ley N.° 19.628 vigente. La Ley N.° 21.719, promulgada en diciembre de 2024, tiene entrada en vigor diferida al 1 de diciembre de 2026; sus exigencias deben incorporarse a la planificación de cumplimiento y confirmarse antes de la operación (\parencite{ley21719}).
  \item \textbf{Marco GLEC (\textit{Global Logistics Emissions Council}) / ISO 14083:2023:} El marco ofrece una metodología para calcular y reportar emisiones de dióxido de carbono equivalente (CO2e) por tonelada-kilómetro transportada (g~CO2e/t-km), considerando las fases \textit{Well-to-Tank} (WTT) y \textit{Tank-to-Wheel} (TTW). El requisito del cliente exportador para 2029 hace necesario identificar las fuentes de datos disponibles y declarar los supuestos utilizados en el cálculo (\parencite{glec2023,iso14083}).
  \item \textbf{Registros de jornada y tacógrafos:} Los registros de tacógrafo y los sistemas de control de asistencia pueden aportar antecedentes sobre conducción y jornada. Su formato, extracción y conservación deben cumplir las reglas aplicables; no se presume que una firma criptográfica por sí sola otorgue valor probatorio ni inalterabilidad jurídica (\parencite{res_ex_1213_2009}).
\end{itemize}

Las referencias oficiales para las leyes citadas se incorporan a la bibliografía compartida. La aplicabilidad jurídica concreta de sus disposiciones a cada situación operativa debe confirmarse antes de la entrega.

\subsection{Modelación cuantitativa de la estacionalidad operativa}\label{sec:2-estacionalidad}

La operación combina variaciones estacionales de demanda con restricciones de conectividad y cierres cordilleranos. La curva estacional de demanda y las restricciones de red deben presentarse con sus fuentes y unidades claramente identificadas.

\marcador{Reemplazar la figura de estacionalidad compartida antes de entregar: la versión disponible afirma un peak de más de 450 viajes diarios sin una base derivable y presenta 288 horas como requisito de retención a bordo; S2 debe mostrar únicamente la autonomía mínima exigida de 72 horas y una curva de demanda sustentada en datos del caso.}

La revisión distingue tres condiciones que inciden en la planificación y que deben validarse con sus fuentes:

\begin{enumerate}
  \item \textbf{Temporada frutícola y de agroexportación (diciembre a abril):} El caso atribuye a estos cinco meses el 65\% de la demanda anual del servicio de frío y reporta esperas de andén que superan las ocho horas, frente a un promedio de 3 horas y 10 minutos. El dato de 96.000 viajes anuales corresponde a la operación total, por lo que no permite derivar directamente un máximo diario para la flota refrigerada. El número de viajes de frío y su curva diaria deben confirmarse antes de presentar un peak cuantitativo. La figura asociada requiere el mismo cotejo.
  \item \textbf{Cierres Climáticos del Paso Fronterizo Los Libertadores (Junio a Septiembre):} El corredor bioceánico de la Ruta 60 CH hacia Mendoza registra aproximadamente 1.900 cruces de camiones al año. El caso reporta cierres climáticos que pueden afectar la conectividad durante períodos prolongados. Para el sistema embarcado, el requisito mínimo de autonomía local es de \textbf{72 horas} (\caso{RT-03.10}{31}); capacidades superiores propuestas para el hardware se describen y justifican en el Subdocumento 4, sin ampliar este requisito base.
  \item \textbf{Otras restricciones operacionales reportadas:} El caso menciona restricciones de circulación durante algunas festividades, con lapsos de 12 a 36 horas, y asignación de ocho analistas al procesamiento mensual de liquidaciones de transportistas. La vigencia, duración y dedicación deben cotejarse con la operación actual.
\end{enumerate}

\clearpage
\section{Dimensionamiento del problema}\label{sec:2-dimensionamiento}

El dimensionamiento cuantitativo del problema operacional de Transportes Curimón S.A. se sustenta en el análisis riguroso de las magnitudes físicas, financieras, territoriales y de seguridad que componen la actividad de la empresa. En la \tabref{2-volumetria} se sintetizan los parámetros volumétricos consolidados que caracterizan la escala de la operación actual y las proyecciones requeridas para la estabilidad del servicio.

\begin{tabla}[ancho=completo]{Volumetría operacional y magnitudes del desafío logístico}{2-volumetria}{L{40mm} C{22mm} C{22mm} L{20mm} X}
\textbf{Métrica Operacional} & \textbf{Línea Base} & \textbf{Proyección 3 Años} & \textbf{Unidad} & \textbf{Impacto en la Operación} \\ \midrule
Viajes Anuales Totales & 96.000 & 118.000 & viajes / año & Promedio de 263 viajes/día; peak estacional pendiente de fuente y base de cálculo. \\
Kilómetros Recorridos & 41.000.000 & 50.000.000 & km / año & Desgaste intensivo de activos y base de odometría para mantenimiento. \\
Kilómetros en Vacío & 10.660.000 & $< 7.500.000$ & km / año & Corresponde al 26\% del kilometraje anual sin generar facturación. \\
Carga Total Transportada & 2.400.000 & 2.900.000 & ton / año & Exposición a 142 detenciones viales anuales por sobrepeso por eje. \\
Documentación D.E.T. & 128.000 & 157.000 & docs / año & Guías y documentos tributarios a emitir y respaldar sin papel. \\
Puntos de Carga y Descarga & 1.400 & 1.700 & recintos & Recintos de clientes donde se producen esperas medias de 3 h 10 min. \\
Vigencias Documentales & 6.000 & 7.000 & fechas & Fechas críticas de conductores y vehículos gestionadas en 4 Excel. \\
Cruces Paso Los Libertadores & 1.900 & 2.400 & cruces / año & Operación binacional sujeta a cierres climáticos. \\
Abastecimientos Diésel & 74.000 & 90.000 & eventos / año & Carga en estanque de San Bernardo y estaciones externas en ruta. \\
Pasadas de Peajes (TAG) & 620.000 & 760.000 & transacciones & Conciliación de peajes interurbanos desfasada sin costeo por ruta. \\
Migración Histórica TMS 2013 & 480.000 & N/A & viajes & Base histórica de 5 años a migrar y conciliar según \caso{RT-05.15}{31}. \\
\end{tabla}

El análisis de la \tabref{2-volumetria} muestra la escala transaccional que debe gestionarse. El 26\% de kilómetros recorridos en vacío (10,66 millones de km) constituye una brecha de utilización de la flota que afecta el resultado operacional.

Las brechas de captura de datos y coordinación pueden afectar el seguimiento de viajes y el análisis de costos. Estas relaciones causales son hipótesis de trabajo que deben contrastarse con datos del mandante.

\marcador{Actualizar la figura causal compartida antes de reincorporarla: su encabezado «Patologías» y algunas relaciones causa-efecto requieren alinearse con el diagnóstico prudente de S2 y cotejarse con datos del mandante.}

\subsection{Análisis de los 7 bloques de datos duros de Curimón}\label{sec:2-datos-duros}

A partir de la cadena causal expuesta, se dimensiona el impacto numérico de los siete bloques de datos duros que estructuran la operación:

\begin{enumerate}
  \item \textbf{Parque Vehicular y Asimetría de Tenencia:} La flota totaliza 374 tractocamiones, divididos en 148 unidades propias (39,6\%, con antigüedad media de 6,4 años) y 226 unidades subcontratadas (60,4\%). A esto se añaden 210 semirremolques propios, de los cuales 44 son furgones refrigerados y 18 unidades están certificadas bajo D.S. N.° 298 para sustancias peligrosas. Los 226 camiones subcontratados pertenecen a 148 transportistas independientes que poseen entre 1 y 4 vehículos. Curimón provee la relación comercial y absorbe la responsabilidad patronal y civil completa, pero carece de potestad sobre el 60,4\% de los tractocamiones motrices, impidiendo cualquier política de estandarización por imposición jerárquica.
  \item \textbf{Fuerza conductora y control de jornada:} El caso considera 454 conductores: 196 dependientes de Curimón y 258 dependientes de 148 transportistas terceros. Informa que no se descargan registros de tacógrafos y describe limitaciones para verificar actividades previas de conductores externos. También reporta cuatro siniestros graves con lesiones atribuidas a somnolencia y fatiga durante los últimos tres años. Estos antecedentes deben cotejarse con los registros correspondientes; la revisión legal debe precisar los controles de jornada exigibles a cada empleador.
  \item \textbf{Red Vial, Kilometraje y Retornos en Vacío:} La flota recorre 41 millones de kilómetros al año a lo largo de un corredor de 3.000 kilómetros. La descoordinación entre la demanda de transporte y la localización de los equipos da lugar a que el 26\% de la distancia total (10,66 millones de km) se recorra en vacío. Esta ineficiencia equivale a movilizar una flota virtual de cerca de 97 tractocamiones consumiendo diésel, peajes y neumáticos sin percibir tarifa alguna.
  \item \textbf{Gestión documental y telemetría:} Cerca de 6.000 vencimientos se administran en cuatro planillas. El caso registra la inmovilización de una unidad SUSPEL durante 14 horas en abril de 2026 por un certificado vencido. Además, 61 tractocamiones propios tienen telemetría CAN de fábrica no integrada, 34 camiones externos carecen de GPS y 192 unidades con GPS se distribuyen entre tres plataformas; dos de ellas no disponen de interfaces API hacia la Torre de Control.
  \item \textbf{Sobreestadías y evidencia de entrega:} Las esperas en los 1.400 recintos de carga y descarga promedian 3 horas y 10 minutos y superan las 8 horas en temporada agrícola. El caso indica que se objeta el 71\% de los cobros por sobreestadía debido a la falta de registros objetivos de llegada, espera y salida. También se informa que el 4,2\% de los comprobantes de entrega en papel se extravía, se vuelve ilegible o se deteriora, lo que demora el ciclo de cobro.
  \item \textbf{Costeo por ruta y contrato:} El caso informa diferencias de rentabilidad entre contratos: tres de los ocho principales operan bajo el costo técnico y concentran el 31\% de los ingresos; uno registra margen negativo de 14\% durante cuatro años. El prorrateo de costos según ingresos dificulta identificar el resultado de cada ruta. La recepción de facturas de combustible puede demorarse hasta 40 días, lo que retrasa el análisis de rendimiento.
  \item \textbf{Seguridad vial y pesajes:} Durante 2025 se registraron 142 detenciones en plazas de pesaje por infracciones de peso por eje, con 2.556 horas-camión de inmovilización. En paralelo, el cliente exportador principal, que representa el 19\% de los ingresos según el caso, condiciona la renovación prevista para 2029 al cumplimiento de requisitos de jornada, documentación digital, posicionamiento y reportes de emisiones GEI bajo GLEC / ISO 14083.
\end{enumerate}

\clearpage
\subsection{Caracterización territorial y mapeo de nodos críticos}\label{sec:2-nodos-criticos}

El despliegue territorial de Transportes Curimón S.A. abarca una red física y logística compleja compuesta por seis tipologías de nodos operacionales, cuya infraestructura actual impone severas restricciones técnicas que deben ser consideradas en el diagnóstico. En la \tabref{2-nodos-criticos} se sintetizan las características y niveles de criticidad de los nodos que conforman la red operacional de la empresa.

\begin{tabla}[ancho=completo]{Síntesis de nodos críticos de la red operacional}{2-nodos-criticos}{L{28mm} L{32mm} L{26mm} X C{16mm}}
\textbf{Nodo Operacional} & \textbf{Tipología e Instalación} & \textbf{Enlaces} & \textbf{Restricción Crítica Identificada} & \textbf{Criticidad} \\ \midrule
Terminal San Bernardo & Casa matriz, torre 24x7, taller central, estanque diésel. & Fibra óptica principal + enlace 4G/5G. & Sala servidores 26 m², split doméstico y UPS 20 min (\caso{RT-06}{32}). & Máxima \\
Terminales Regionales (4) & Valparaíso, Concepción, Antofagasta, Puerto Montt. & Enlaces comerciales locales; 3 de 4 sin respaldo. & Vulnerabilidad de enlace; puntos de relevo macrozonas norte, centro y sur. & Alta \\
Talleres Propios (2) & San Bernardo y Los Ángeles; 46 operarios técnicos. & Conectados a red de base; apoyo en Talca. & Ciclo de paso cada 6 días en propios; terceros pasan cada 30 días o más. & Alta \\
Paso Los Libertadores & Corredor binacional a Mendoza ($\approx 1.900$ viajes). & Conectividad celular precaria de montaña. & Cierres climáticos por nieve que interrumpen la señal. & Crítica \\
Zonas de Sombra Celular & Ruta 5 Norte (Atacama) y tramos cordilleranos. & Cero cobertura en tramos $> 80$ km. & Pérdida de transmisión telemática en tiempo real y bloqueo e-Docs. & Crítica \\
Puntos de Clientes ($\approx 1.400$) & Plantas, packings, mineras y recintos portuarios. & Infraestructura ajena; sin equipos fijos. & Esperas no acreditadas (3 h 10 min a $> 8$ h); 71\% sobreestadías objetadas. & Alta \\
\end{tabla}

El análisis de la \tabref{2-nodos-criticos} muestra que la sala de servidores del Terminal San Bernardo (26 m², climatización tipo split y UPS con 20 minutos de autonomía) no satisface los requisitos de infraestructura citados en las Bases Técnicas Transversales (\caso{RT-06.01}{32} a \caso{RT-06.09}{32}). Por ello, el uso de estas dependencias para servicios de operación continua requiere evaluar las brechas y medidas de continuidad correspondientes.

Asimismo, la presencia de tramos con más de 80 kilómetros continuos sin cobertura de telecomunicaciones en la Ruta 5 Norte y los cierres cordilleranos reportados en el Paso Los Libertadores determinan que la arquitectura tecnológica no puede asumir la conectividad permanente como un supuesto válido. La desconexión es una condición física intrínseca a la geografía chilena; conforme al requisito base, los sistemas embarcados deben retener información localmente por al menos 72 horas y sincronizarla al recuperar señal celular.

\section{Actores y grupos de interés}\label{sec:2-actores}

La operación y gobernanza de Transportes Curimón S.A. involucra a un ecosistema diverso de trece (13) actores clave, cuyos intereses, expectativas y capacidades de bloqueo determinan la viabilidad de cualquier transformación en los procesos corporativos. En el diagnóstico previo se omitió a tres actores estratégicos: el \textbf{Fondo de Inversión Institucional}, la \textbf{Dirección del Trabajo (DT)} y la \textbf{Aseguradora de Carga y Flota}.

En la \figref{2-matriz-actores} se ilustra el mapeo de actores en la Matriz de Poder e Influencia frente al Nivel de Interés, identificando la estrategia de gestión corporativa requerida para cada uno.

\figura{figuras/02-problema/Matriz_Interes_Poder_Stakeholders.png}{Matriz de poder frente a interés y mapa de influencia de los 13 actores}{2-matriz-actores}

El análisis de la \figref{2-matriz-actores} permite estructurar la gobernanza de los grupos de interés en cuatro cuadrantes de acción bien diferenciados:
\begin{itemize}
  \item \textbf{Cuadrante 1 — Gestionar de Cerca (Alto Poder / Alto Interés):} Concentra a las gerencias de Curimón (General, Operaciones, Finanzas y Prevención) y al Cliente Exportador del 19\%. Estos actores son los garantes de la viabilidad económica y la continuidad del negocio; una fricción no resuelta con cualquiera de ellos paraliza la operación o deriva en la pérdida del contrato principal.
  \item \textbf{Cuadrante 2 — Mantener Satisfecho (Alto Poder / Bajo Interés Diario):} Incluye a las entidades fiscalizadoras y de capital: la Dirección del Trabajo (DT), la Aseguradora de Carga, el Fondo de Inversión Institucional (22\%) y la Familia Fundadora (78\%). Estos actores no intervienen en el despacho diario, pero poseen facultad de clausura legal, revocación de pólizas de seguro, o veto financiero sobre el Directorio.
  \item \textbf{Cuadrante 3 — Monitorear (Bajo Poder / Bajo Interés Operativo Central):} Agrupa a las jefaturas técnicas intermedias (TI y Mantenimiento), que canalizan la viabilidad instrumental de las plataformas.
  \item \textbf{Cuadrante 4 — Mantener Informado y Asegurar Adhesión (Bajo Poder Jerárquico / Alto Interés de Campo):} Compuesto por los 148 dueños subcontratados, los 258 choferes externos y los 196 choferes propios. Aunque individualmente carecen de poder societario, su poder colectivo de veto de facto es altísimo: si los transportistas rechazan compartir datos o los conductores sabotean el registro de jornada, la operación de Curimón queda desabastecida en un 60,4\%.
\end{itemize}

En la \tabref{2-actores-sintesis} se sintetizan las características, expectativas y riesgos asociados a cada uno de los trece actores del ecosistema.

\begin{tabla}[ancho=completo]{Matriz sintética de actores y grupos de interés}{2-actores-sintesis}{L{30mm} L{26mm} X L{24mm} L{24mm}}
\textbf{Actor} & \textbf{Dotación / Poder} & \textbf{Dolor Operacional Principal} & \textbf{Capacidad Bloqueo} & \textbf{Nivel de Riesgo} \\ \midrule
1. Fondo Inversión & 22\% propiedad. & Pasivos laborales contingentes y caída del margen (9,0\%). & Máxima (Directorio) & Financiero / Gobierno \\
2. Dirección Trabajo & Ente fiscalizador. & Cero descargas de tacógrafos e infracción Art. 25 bis. & Extrema (Clausura) & Legal / Regulatorio \\
3. Aseguradora Carga & Cobertura pólizas. & Opacidad en siniestros y falta de registro térmico. & Alta (Cobertura) & Financiero / Pólizas \\
4. Familia Fundadora & 78\% propiedad. & Exposición por la renovación del contrato principal en 2029. & Alta (Societaria) & Estratégico \\
5. Gerencia General & E. Valdebenito. & Fractura entre responsabilidad 100\% y control terceros. & Máxima (Ejecutiva) & Operacional y penal \\
6. Gerencia Operaciones & R. Mansilla. & Asignación a ciegas, 26\% retornos vacíos y 3 GPS. & Alta (Despacho) & Fricción operacional \\
7. Gerencia Finanzas & G. Ossandón. & 3 contratos a pérdida (-14\%), diésel con 40 d desfase. & Alta (Financiera) & Rentabilidad \\
8. Mantenimiento & H. Trincado. & Preventivo por adivinanza; 61 CAN de fábrica inactivos. & Media (Logística) & Falla en ruta \\
9. Prevención Riesgos & D. Aguayo. & Responsabilidad solidaria; 6.000 vigencias en Excel. & Alta (Veto legal) & Siniestralidad / Multas \\
10. TI y Flota & M. Riquelme / P. Kast. & TMS 2013 rígido, 34 camiones sin GPS y 9 personas TI. & Alta (Técnica) & Obsolescencia / Silos \\
11. Choferes Propios & 196 choferes. & Fatiga en ruta, bermas inseguras, esperas >8 h andén. & Alta (De facto) & Seguridad / Sindical \\
12. Terceros Externos & 148 dueños. & Vulneración soberanía activo, liquidación lenta y errores. & Muy Alta (Colectiva) & Desabastecimiento 60\% \\
13. Cliente Exportador & A. Lecaros (19\%). & Incumplimiento exigencias 2029 (CO2e GLEC, e-Docs). & Extrema (Comercial) & Quiebre facturación \\
\end{tabla}

\subsection{Principios de arbitraje operacional de las seis tensiones estructurales}\label{sec:2-arbitraje}

La coexistencia de estos trece actores plantea seis tensiones de coordinación y gobernanza que se resumen a continuación:

\begin{enumerate}
    \item \textbf{Privacidad de conductores y transportistas externos frente a la necesidad de control operacional:} El caso plantea la necesidad de verificar la jornada y la ubicación durante los viajes asignados, junto con la preocupación de transportistas externos por el seguimiento fuera de esos períodos. La Ley N.° 21.719 entra en vigor el 1 de diciembre de 2026; hasta entonces deben considerarse las disposiciones vigentes de la Ley N.° 19.628. La base jurídica, finalidad, información a titulares y límites de tratamiento requieren revisión legal.
  \begin{itemize}
    \item \textit{Criterio de diseño por validar:} Limitar la captura a los datos necesarios para la finalidad operacional acordada y definir períodos de retención, acceso y eliminación. La base jurídica y los mecanismos de información o autorización deben confirmarse antes de procesar datos personales; no se presume que el consentimiento sea la única base aplicable.
  \end{itemize}
  \item \textbf{Flexibilidad de Despacho Manual frente a Asignación Rigurosa y Seguridad Vial:} Operaciones busca preservar la discrecionalidad de los 22 despachadores para autorizar salidas con documentación en trámite y así cumplir itinerarios comerciales, mientras que Prevención de Riesgos exige el bloqueo estricto ante cualquier vencimiento de vigencias o límites de jornada.
  \begin{itemize}
    \item \textit{Criterio de diseño por validar:} Definir con Operaciones, Prevención y asesoría laboral qué documentos y verificaciones deben bloquear una asignación. Los casos de excepción, responsables y registros de auditoría deben acordarse con el mandante y las bases antes de implementarse.
  \end{itemize}
  \item \textbf{Autonomía del Transportista Tercero frente a Estandarización de Datos de Flota:} Curimón requiere visibilidad sobre el 60,4\% de la flota subcontratada, pero carece de potestad jurídica para forzar a 148 empresarios independientes a sustituir sus sistemas GPS actuales o a realizar inversiones obligatorias de modernización.
  \begin{itemize}
    \item \textit{Criterio de Arbitraje y Principio Rector:} Integración no traumática basada en homologación progresiva e incentivos. El modelo de gestión debe admitir la heterogeneidad de fuentes preexistentes mediante interfaces estandarizadas de datos, focalizando la provisión de nuevo equipamiento exclusivamente en las unidades desprovistas de seguimiento, y recompensando la entrega fidedigna de información operativa con transparencia en las liquidaciones mensuales de fletes.
  \end{itemize}
  \item \textbf{Presión Comercial de Entrega Inmediata frente a Restricciones Operativas de Descanso:} La fuerza de ventas y los clientes exigen tiempos de tránsito acelerados para cumplir ventanas de descarga portuaria o faenas mineras, induciendo indirectamente a los choferes a exceder las 5 horas continuas de conducción (Art. 25 bis del Código del Trabajo).
  \begin{itemize}
    \item \textit{Criterio de diseño:} Estimar tiempos de viaje considerando las restricciones de conducción y descanso aplicables, condiciones de ruta y puntos seguros de detención. Los parámetros y fuentes deben validarse antes de convertirlos en compromisos de planificación.
  \end{itemize}
  \item \textbf{Costeo por ruta y contrato:} La asignación histórica de costos en función de los ingresos dificulta identificar las pérdidas de los contratos que operan bajo costo y comparar el resultado de cada ruta.
  \begin{itemize}
    \item \textit{Criterio de análisis:} Comparar alternativas de imputación de costos directos (combustible, peajes y fletes a terceros) a cada viaje y contrato. El método debe definirse con Finanzas y validarse frente a las fuentes contables disponibles.
  \end{itemize}
  \item \textbf{Exigencia de Trazabilidad Integral del Cliente Exportador frente a Heterogeneidad Tecnológica:} El cliente principal (19\% de la facturación) exige un estándar unificado de datos para la renovación contractual de 2029, mientras que Curimón opera con un parque mixto compuesto por camiones propios y de terceros con dispares niveles de sensorización y plataformas aisladas.
  \begin{itemize}
    \item \textit{Criterio de análisis:} Considerar la normalización de datos de distintas fuentes para facilitar la reportabilidad. La cobertura, interoperabilidad y calidad de datos deben verificarse antes de comprometer resultados independientes de los dispositivos en uso.
  \end{itemize}
\end{enumerate}

\section{Resumen de Requerimientos, Supuestos, Exclusiones y Restricciones}\label{sec:2-requerimientos}

El cierre analítico del diagnóstico operacional consolida las condiciones de contorno que delimitan el alcance del proyecto. Conforme a las directrices de ingeniería, en este acápite se sintetizan las necesidades del mandante y se formula una matriz de supuestos de ingeniería que evalúa el impacto de eventuales contingencias.

\subsection{Síntesis de requerimientos del negocio y operacionales preliminares}\label{sec:2-req-sintesis}

Las necesidades levantadas a partir de las Bases Técnicas del Caso 10 y las sesiones de trabajo con las gerencias de Curimón se agrupan en cuatro dominios funcionales clave:
\begin{enumerate}
  \item \textbf{Dominio de asignación y control operacional:} Analizar la factibilidad de consultar la disponibilidad del conductor, idoneidad del vehículo, vigencias documentales y compatibilidad de carga antes de asignar un viaje. El tiempo objetivo de respuesta y las verificaciones bloqueantes deben definirse con la RTM, las bases y el mandante.
  \item \textbf{Dominio de trazabilidad y gestión de esperas:} Evaluar el registro de arribos, esperas y salidas en los recintos pertinentes para aportar antecedentes a la revisión de cobros por sobreestadía. La cobertura, precisión y valor probatorio de los registros deben acordarse y probarse.
  \item \textbf{Dominio de documentación digital y continuidad en zonas sin cobertura:} Evaluar la emisión y sincronización diferida de los documentos considerados en el caso (aproximadamente 128.000 al año), su integración con sistemas externos y los controles requeridos para operar sin conectividad. Los requisitos legales y tributarios deben confirmarse antes de definir una emisión sin conexión.
  \item \textbf{Dominio de costeo y sostenibilidad:} Evaluar el cálculo del costo directo por viaje y de emisiones por tonelada-kilómetro bajo el marco GLEC / ISO 14083. La disponibilidad y calidad de las fuentes de datos, periodicidad y cobertura de flota deben confirmarse con el mandante.
\end{enumerate}

En el plano de los requerimientos no funcionales y de resiliencia, se resumen las exigencias aplicables de las bases y el requisito de operación desconectada:
\begin{itemize}
  \item \textbf{Disponibilidad de los servicios críticos:} $\ge 99,9\%$, conforme al \art{20}{14} y a \transv{RT-10.01}{22}.
  \item \textbf{Tiempo de Recuperación (RTO):} $\text{RTO} \le 4\text{ horas}$, conforme al \art{20}{14} y a \transv{RT-07.04}{17}.
  \item \textbf{Punto de Recuperación (RPO):} $\text{RPO} \le 15\text{ minutos}$, conforme al \art{20}{14} y a \transv{RT-07.04}{17}.
  \item \textbf{Autonomía telemática desconectada:} Al menos 72 horas de retención local, conforme a \caso{RT-03.10}{31}; la capacidad efectiva requiere prueba sobre el hardware seleccionado.
\end{itemize}

\subsection{Matriz de supuestos auténticos de ingeniería de proyectos}\label{sec:2-supuestos-matriz}

La \tabref{2-supuestos} organiza supuestos preliminares de planificación, con una estimación inicial de probabilidad e impacto y posibles medidas de respuesta. Estas valoraciones no constituyen un análisis de riesgos validado; deben revisarse con los responsables de cada dependencia y coordinarse con el Subdocumento 8.

\begin{tabla}[ancho=completo]{Matriz de supuestos de ingeniería del proyecto}{2-supuestos}{L{44mm} C{16mm} C{18mm} X}
\textbf{Supuesto de Ingeniería} & \textbf{Prob.} & \textbf{Impacto} & \textbf{Estrategia de Mitigación y Contingencia Técnica} \\ \midrule
SUP-01: Adhesión de Terceros (escenario objetivo 85\%) & Media (estimación) & Alto & Si el 85\% de los 226 tractocamiones de terceros comparte datos, unos 34 quedarían fuera de ese nivel de integración. Confirmar la meta con el mandante y evaluar enrolamiento e incentivos. \\
SUP-02: Acceso a Interfaces de Terceros & Media (estimación) & Alto & Si no se habilitan interfaces para las tres plataformas existentes, hasta 192 unidades con GPS quedarían sin vista unificada. Confirmar permisos, cobertura y condiciones de servicio; evaluar alternativas de integración u homologación. \\
SUP-03: Continuidad de Servicios Públicos (DT y SII) & Baja (estimación) & Alto & Una interrupción podría retrasar verificaciones laborales o trámites documentales. Definir dependencias y contingencia con asesoría normativa y técnica; dimensionar el impacto con datos de disponibilidad. \\
SUP-04: Resiliencia en Cordillera (cierres prolongados reportados) & Baja (estimación) & Alto & Asegurar al menos 72 h de retención local conforme a RT-03.10; cualquier margen adicional requiere dimensionamiento con el perfil real de eventos y pruebas de desconexión y recuperación. \\
SUP-05: Integridad de Garantías Vehiculares & Media (estimación) & Alto & Una incompatibilidad del método de captura podría afectar hasta 61 tractocamiones propios con CAN J1939 de fábrica. Confirmar condiciones de garantía e interfaces con fabricantes y probar el método seleccionado. \\
SUP-06: Cadencia de Ingreso a Terminales & Media (estimación) & Alto & Una demora en accesos podría afectar la instalación propuesta de hasta 121 nuevas pasarelas. Confirmar ventanas de acceso y mantenimiento con cada terminal y acordar un plan con Operaciones. \\
\end{tabla}

\subsection{Exclusiones y restricciones principales}\label{sec:2-exclusiones}

Para fijar la frontera formal del proyecto y evitar desviaciones de alcance, se declaran las siguientes exclusiones y restricciones contractuales:
\begin{itemize}
  \item \textbf{Exclusiones Contractuales Explícitas:}
  \begin{enumerate}
    \item No forman parte del alcance la provisión de combustible físico, repuestos mecánicos ni el mantenimiento preventivo/correctivo de los motores y carrocerías de los vehículos.
    \item No se incluye el licenciamiento, modificación del código fuente base ni la administración del ERP transaccional contable de Curimón, interactuando con este exclusivamente a través de interfaces API estandarizadas.
    \item Se excluyen obras civiles mayores de remodelación edilicia en las salas de servidores de terminales, orientándose la arquitectura hacia servicios en la nube pública de alta resiliencia.
  \end{enumerate}
  \item \textbf{Restricciones Operacionales y de Entorno:}
  \begin{enumerate}
    \item \textbf{Seguridad de interacción en cabina:} Como condición de diseño propuesta, reducir las interacciones que aparten la atención del conductor durante la marcha. El alcance de bloqueos, alertas y excepciones debe revisarse conforme a la Ley N.° 21.377 y la reglamentación aplicable.
    \item \textbf{Restricción de Cadena de Frío:} Registro térmico ininterrumpido en el rango de -30 °C a +30 °C con resolución de 0,1 °C para las 44 ramplas refrigeradas.
    \item \textbf{Restricción sobre montos de la oferta (\art{50.2}{32}):} El documento técnico no incorpora tarifas, honorarios de desarrollo ni costos de la oferta.
  \end{enumerate}
\end{itemize}

\vspace{0.4cm}
\noindent Los catálogos exhaustivos de requerimientos preliminares, el inventario detallado del parque vehicular y conductores, la matriz exhaustiva de restricciones legales y las trece fichas completas de caracterización de actores se entregan en el documento independiente \texttt{AUDIT-Subdocumento2-Anexos.pdf}, el cual complementa el presente diagnóstico conforme a las directrices del Comunicado 10.

\nocite{bases-fep01, bases-fep02, bases-fep03, bases-com10, ley20123, codigo_trabajo_25bis, ley21719, ley21377, glec2023, iso14083, ds298, ds43minsal, ds158mop}
